\documentclass[a4paper,11pt]{ctexart}
\usepackage{xcolor}
\usepackage{array}
\usepackage{booktabs}
\usepackage{hyperref}
\hypersetup{colorlinks=true,linkcolor=blue!70!black}
\linespread{1.35}

\newcommand{\draft}{\textcolor{orange!80!black}{\textbf{【待写】}}}

\title{\textcolor{blue!70!black}{\Huge\textbf{法学模块 · 课程大纲}}}
\author{}
\date{}

\begin{document}
\maketitle
\thispagestyle{empty}

\section*{第零层 · 入门故事}

\begin{enumerate}
\item[\textbf{第1课}] \textbf{为什么要有红绿灯？——规则如何保护所有人} \draft \\
\textbf{内容：} 没有规则的路口=混乱。法律的作用：让人们的行为可预测。从红绿灯到交通法到刑法的递进。
\item[\textbf{第2课}] \textbf{包青天的故事——古代法官怎么判案} \draft \\
\textbf{内容：} 包拯审案的方法（证据/推理/证人）。古代的衙门制度。对比今天的法庭——法官/检察官/律师/被告各司其职。
\item[\textbf{第3课}] \textbf{十二怒汉——陪审团是怎么工作的？} \draft \\
\textbf{内容：} 1957年经典电影《十二怒汉》——一个男孩被指控杀人，11个陪审员认为有罪，1个认为存疑。一个人的坚持如何改变了整个裁决。为什么"合理怀疑"是刑事审判的基石？
\end{enumerate}

\section*{深入课程}

\begin{enumerate}
\item[\textbf{第1课}] \textbf{法理学——什么是"法"？} \draft \\
\textbf{内容：} 自然法 vs 实证法（Legal Positivism）的千年争论。法律与道德的关系。法的效力层级（宪法→法律→法规→规章）。
\item[\textbf{第2课}] \textbf{宪法与行政法——权力怎么被关进笼子} \draft \\
\textbf{内容：} 宪法的核心：限制公权力、保障私权利。三权分立、违宪审查制度。中国宪法→基本法→行政法的结构。情景：齐玉苓案（中国宪法司法化第一案）为什么重要？
\item[\textbf{第3课}] \textbf{刑法——犯罪与刑罚的边界} \draft \\
\textbf{内容：} 罪刑法定（法无明文规定不为罪）、四要件犯罪构成（客体/客观方面/主体/主观方面）。正当防卫 vs 防卫过当。情景：昆山龙哥反杀案（2018）——为什么检察院认定是正当防卫？
\item[\textbf{第4课}] \textbf{民法——每天都在签的合同} \draft \\
\textbf{内容：} 民事主体、合同成立与生效的四个要件（意思表示真实/主体适格/内容合法/形式符合）、合同违约与救济。侵权（过错/无过错/公平责任）。情景：买了一杯奶茶发现里面有虫子——消费者权益保护怎么维权？
\item[\textbf{第5课}] \textbf{刑事诉讼法——程序正义为什么和实体正义一样重要？} \draft \\
\textbf{内容：} 无罪推定、非法证据排除、"毒树之果"理论。侦查→审查起诉→审判→执行。情景：聂树斌案（1995年被执行死刑，2016年改判无罪）——冤案是怎么发生的？纠正冤案需要什么制度保障？
\item[\textbf{第6课}] \textbf{北京在地——北京的法院与法律资源} \draft \\
\textbf{内容：} 最高人民法院（位于南四环）、北京市高级人民法院、海淀法院（全国收案数量最大的基层法院之一）。北京的法律服务业——国贸CBD的外所与中国红圈所。
\end{enumerate}

\end{document}
